% ======================================================================
%  sec27.tex —— 第 27 节 相对 Proj
%  底本：原书 pp.133–140；原书标题：27. Relative Proj
% ======================================================================

% ---------- 本节用到的宏（若全书导言区已定义则自动跳过） ----------
\providecommand{\QCoh}{\mathsf{QCoh}}
\providecommand{\Sym}{\operatorname{Sym}}

\section{相对 $\Proj$}\label{sec:27}

本节固定一个概形 $Y$，以及一个 $\Z_{\ge 0}$-分次拟凝聚 $\cO_Y$-代数 $S$。按定义，这意味着有一个拟凝聚 $\cO_Y$-代数 $S$，配有 $S$ 作为 $\cO_Y$-模的分解
\[
S = \bigoplus_{n \in \Z_{\ge 0}} S_n,
\]
且该分解与 $S$ 上的乘法结构相容。

\subsection{相对 Proj 构造}\label{sec:27.1}

下述构造把 $\Proj(S)$ 推广到相对情形。

\begin{construction}\label{con:27.1.1}
对任意仿射开子集 $U \subseteq Y$，限制 $S|_U$ 对应于 $\Z_{\ge 0}$-分次 $\cO_Y(U)$-代数 $S(U)$。考虑概形 $\Proj(S(U))$ 与典范态射
\[
p_U \colon \Proj(S(U)) \to \Spec(S(U)) \simeq U.
\]
对仿射开子集 $U \subseteq V \subseteq Y$，因 $S$ 拟凝聚，有
\[
S(U) \simeq S(V) \otimes_{\cO_Y(V)} \cO_Y(U).
\]
于是由引理 \ref{lem:25.3.5}，有典范的笛卡儿方块
\[
\begin{tikzcd}[column sep=3em, row sep=2.6em]
\Proj(S(U)) \arrow[r] \arrow[d, "p_U"'] & \Proj(S(V)) \arrow[d, "p_V"] \\
U \arrow[r] & V.
\end{tikzcd}
\]
由此可见，存在本质唯一的 $Y$-概形，记作
\[
p \colon \Proj_Y(S) \to Y,
\]
并对任意仿射开子集 $U \subseteq Y$ 配有 $U$-同构
\[
\Proj_Y(S)_U \simeq \Proj(S(U)),
\]
使得对仿射开子集 $U \subseteq V \subseteq Y$，下图交换：
\[
\begin{tikzcd}[column sep=3.4em, row sep=2.6em]
\Proj_Y(S)_U \arrow[r] \arrow[d, "\simeq"'] & \Proj_Y(S)_V \arrow[d, "\simeq"] \\
\Proj(S(U)) \arrow[r] & \Proj(S(V)).
\end{tikzcd}
\]
\end{construction}

态射 $p \colon \Proj_Y(S) \to Y$ 分离，因为每个 $\Proj_Y(S)_U$ 都分离。

\begin{yizhu}
构造 \ref{con:27.1.1} 把相对 $\Proj$ 定义为局部构造的粘接：在每块仿射开集上取 $\Proj(S(U))$，再用引理 \ref{lem:25.3.5} 的笛卡儿性保证粘接数据相容。Stacks Project 的相应陈述要求 $S$ 是拟凝聚分次 $\cO_Y$-代数，并把粘接同构一并记下：对仿射开集 $U \subseteq U'$，除概形的同构外还有 $\theta \colon r^*\cO_{\Proj(A')}(n) \simeq \cO_{\Proj(A)}(n)$。扭转层必须跟着一起粘，否则 $\cO(1)$ 只是逐块的层而非 $\Proj_Y(S)$ 上的整体对象。
[来源：Stacks Project, Tag 01NM（Situation 27.15.1 与 Lemma 27.15.2，Relative Proj via glueing），\url{https://stacks.math.columbia.edu/tag/01NM}（2026-11-17 访问）。]
\end{yizhu}

\begin{lemma}\label{lem:27.1.2}
设 $f \colon Y' \to Y$ 为概形态射。则有典范同构
\[
\Proj_Y(S) \times_Y Y' \simeq \Proj_{Y'}(f^*S).
\]
\end{lemma}

\begin{proof}
由引理 \ref{lem:25.3.5} 得到。
\end{proof}

下述构造把 $\Proj(S)$ 的标准开子集推广到相对情形。

\begin{construction}\label{con:27.1.3}
对任意 $f \in S_n(Y)$（$n > 0$），存在唯一的开子集 $U_+(f) \subseteq \Proj_Y(S)$，使得对任意开子集 $V \subseteq Y$，有
\[
U_+(f)_V \simeq U_+(f|_V)
\]
（作为 $\Proj_Y(S)_V \simeq \Proj(S(V))$ 的开子集）。

由构造，$U_+(f) \to Y$ 仿射，且
\[
U_+(f) \simeq \Spec_Y\Bigl(\bigl(\bigoplus_{d \in \Z_{\ge 0}} S_{nd}\bigr)/(f-1)\Bigr).
\]
\end{construction}

下述构造把 $\Proj(S)$ 关于 $S$ 的函子性推广到相对情形。

\begin{construction}\label{con:27.1.4}
设 $\phi \colon S \to S'$ 是 $\Z_{\ge 0}$-分次拟凝聚 $\cO_Y$-代数之间的分次同态。

存在唯一的开子集 $U_+(\phi) \subseteq \Proj_Y(S')$，使得对任意开子集 $V \subseteq Y$，有
\[
U_+(\phi)_V \simeq U_+(\phi|_V)
\]
（作为 $\Proj_Y(S')_V \simeq \Proj(S'(V))$ 的开子集）。

存在唯一的态射 $U_+(\phi) \to \Proj_Y(S)$，使得它在任意开子集 $V \subseteq Y$ 上的限制，经由上一段的等同，等于典范态射 $U_+(\phi)_V \to \Proj(S(V))$。
\end{construction}

\subsection{拟凝聚模}\label{sec:27.2}

下述构造把 $\widetilde{M}$ 推广到相对情形。

\begin{construction}\label{con:27.2.1}
设 $M$ 为 $\Z$-分次拟凝聚 $S$-模。存在唯一的对象
\[
\widetilde{M} \in \QCoh(\Proj_Y(S)),
\]
并对任意仿射开子集 $U \subseteq Y$ 配有同构
\[
\widetilde{M}|_{\Proj(S(U))} \simeq \widetilde{M(U)} \in \QCoh(\Proj(S(U))),
\]
使得这些同构对 $U$ 具有函子性。
\end{construction}

我们有函子
\begin{equation}\label{eq:27.1}
S\text{-}\mathsf{mod}_{\mathrm{gr},\mathrm{qcoh}} \to \QCoh(\Proj_Y(S)), \qquad M \mapsto \widetilde{M},
\end{equation}
它正合且与小余极限交换。

\begin{definition}\label{def:27.2.2}
称增广理想 $S_+$ 由 $S_1$ 生成，若存在由仿射开子集组成的覆盖 $Y = \bigcup_{i \in I} U_i$，使得对每个 $i \in I$，增广理想 $S_+(U_i)$ 由 $S_1(U_i)$ 生成。
\end{definition}

\begin{remark}\label{rem:27.2.3}
容易看出：$S_+$ 作为 $S$ 的理想由 $S_1$ 生成，当且仅当 $S$ 作为 $S_0$-代数由 $S_1$ 生成。
\end{remark}

\begin{corollary}\label{cor:27.2.4}
设 $S_+$ 由 $S_1$ 生成。则
\[
\cO_{\Proj_Y(S)}(1) := \widetilde{S(1)} \in \QCoh(\Proj_Y(S))
\]
是 $\Proj_Y(S)$ 上的线丛。
\end{corollary}

\begin{proof}
由命题—定义 \ref{prop:26.2.2} 得到。
\end{proof}

下述构造把 $\Gamma(\cF)_{\mathrm{gr}}$ 推广到相对情形。

\begin{construction}\label{con:27.2.5}
设 $S_+$ 由 $S_1$ 生成。

函子 (27.1) 有右伴随，由
\begin{equation}\label{eq:27.2}
\QCoh(\Proj_Y(S)) \to S\text{-}\mathsf{mod}_{\mathrm{gr},\mathrm{qcoh}}, \qquad \cF \mapsto p_*(\cF)_{\mathrm{gr}} := \bigoplus_{n \in \Z} p_*\bigl(\cF(n)\bigr)
\end{equation}
给出。
\end{construction}

\subsection{有限型情形}\label{sec:27.3}

\begin{corollary}\label{cor:27.3.1}
设 $S_+$ 由 $S_1$ 生成，且 $S_1$ 有限型。则态射 $\Proj_Y(S) \to Y$ 有限型。
\end{corollary}

\begin{proof}
由引理 \ref{lem:25.4.6} 得到。
\end{proof}

\begin{definition}\label{def:27.3.2}
设 $M$ 为 $\Z$-分次拟凝聚 $S$-模。
\begin{itemize}[leftmargin=2em,itemsep=0.25ex]
  \item 称 $M$ \emph{局部最终有限型}，若存在 $Y$ 的仿射开覆盖 $Y = \bigcup U_i$，使对每个 $i$ 存在 $N_i$，使得
  \[
  \bigoplus_{n \ge N_i} M_n(U_i)
  \]
  是有限生成 $S(U_i)$-模。
  \item 称 $M$ \emph{局部最终为零}，若存在 $Y$ 的仿射开覆盖 $Y = \bigcup U_i$，使对每个 $i$ 存在 $N_i$，使得当 $n \ge N_i$ 时 $M_n(U_i) = 0$。
\end{itemize}
\end{definition}

\begin{corollary}\label{cor:27.3.3}
设 $S_+$ 由 $S_1$ 生成，且 $S_1$ 有限型。
\begin{enumerate}[label=(\roman*),leftmargin=2.6em,itemsep=0.25ex]
  \item 设 $M$ 为局部最终有限型的 $\Z$-分次拟凝聚 $S$-模。则 $\widetilde{M}$ 有限型。
  \item 在上述假设下，$\widetilde{M} = 0$ 当且仅当 $M$ 局部最终为零。
\end{enumerate}
\end{corollary}

\begin{proof}
(i) 由引理 \ref{lem:26.4.2} 得到；(ii) 由习题 \ref{exc:26.4.3} 得到。
\end{proof}

\begin{corollary}\label{cor:27.3.4}
设 $S_+$ 由 $S_1$ 生成，且 $S_1$ 有限型。则函子 (27.2) 全忠实。
\end{corollary}

\begin{proof}
由定理 \ref{thm:26.4.6} 得到。
\end{proof}

\begin{corollary}\label{cor:27.3.5}
设 $S_+$ 由 $S_1$ 生成，且 $S_1$ 有限型。则任意对象 $\cF \in \QCoh(\Proj_Y(S))$ 都形如 $\widetilde{M}$，其中 $M \in S\text{-}\mathsf{mod}_{\mathrm{gr},\mathrm{qcoh}}$。此外，若还有：
\begin{itemize}[leftmargin=2em,itemsep=0.25ex]
  \item $Y$ 拟紧且拟分离；
  \item $\cF$ 有限型，
\end{itemize}
则可取 $M$ 有限型。
\end{corollary}

\begin{proof}
只有最后一个断言需要证明。取任意拟凝聚 $\Z$-分次 $S$-模 $M$ 使 $\cF \simeq \widetilde{M}$。由定理 18.6.6，可把 $M$ 写成滤过余极限
\[
M \simeq \colim_{\alpha} M_\alpha,
\]
其中每个 $M_\alpha \subseteq M$ 是有限型的拟凝聚 $\Z$-分次 $S$-子模。有
\[
\cF \simeq \widetilde{M} \simeq \colim_{\alpha} \widetilde{M_\alpha}.
\]
因 $\cF$ 有限型且 $\Proj_Y(S)$ 拟紧，故当指标 $\alpha$ 充分大时 $\cF \simeq \widetilde{M_\alpha}$。
\end{proof}

\begin{corollary}\label{cor:27.3.6}
设：
\begin{itemize}[leftmargin=2em,itemsep=0.25ex]
  \item 概形 $Y$ 拟紧；
  \item 理想 $S_+$ 由 $S_1$ 生成，且 $S_1$ 有限型。
\end{itemize}
则对任意有限型的 $\cF \in \QCoh(\Proj_Y(S))$，典范态射
\[
p^* p_*\cF(n) \to \cF(n)
\]
当 $n \gg 0$ 时是满射。此外，若 $Y$ 还拟分离，则当 $n \gg 0$ 时存在满射
\[
p^* N \to \cF(n),
\]
其中 $N$ 是有限型的拟凝聚 $\cO_Y$-模。
\end{corollary}

\begin{proof}
第一个断言由命题 \ref{prop:26.4.8} 得到。第二个断言由第一个断言，沿用推论 \ref{cor:27.3.5} 证明中的论证得到。
\end{proof}

在 $\cF$ 为结构层的特殊情形，我们有更强的结果：

\begin{corollary}\label{cor:27.3.7}
设 $S_+$ 由 $S_1$ 生成。则对任意 $n \ge 0$，典范态射
\[
S_n \to p_*\bigl(\cO_{\Proj_Y(S)}(n)\bigr)
\]
诱导满射
\[
\phi_n \colon p^* (S_n) \to \cO_{\Proj_Y(S)}(n).
\]
\end{corollary}

\begin{proof}
由习题 \ref{exc:26.4.9} 得到。
\end{proof}

\begin{remark}\label{rem:27.3.8}
注意 $\phi_0$ 显然是满射；而 $\phi_1$ 是满射蕴涵对任意 $n > 0$，$\phi_n$ 也是满射。
\end{remark}

\begin{corollary}\label{cor:27.3.9}
对任意分次理想 $I \subseteq S$，同态 $S \to S/I$ 诱导闭浸入
\begin{equation}\label{eq:27.3}
\Proj_Y(S/I) \to \Proj_Y(S),
\end{equation}
其定义理想由 $\widetilde{I}$ 给出。此外：
\begin{enumerate}[label=(\arabic*),leftmargin=2.4em,itemsep=0.25ex]
  \item 若 $S_+$ 由 $S_1$ 生成，则沿 (27.3) 的拉回函子把 $\cO(n)$ 送到 $\cO(n)$。
  \item 若 $S_+$ 由 $S_1$ 生成且 $S_1$ 有限型，则任意闭浸入 $i \colon X \to \Proj_Y(S)$ 都形如 (27.3)，其中
  \[
  I := \ker\bigl(S \to p_*(\cO_{\Proj_Y(S)})_{\mathrm{gr}} \to p_*(i_*\cO_X)_{\mathrm{gr}}\bigr).
  \]
\end{enumerate}
\end{corollary}

\begin{proof}
由定理 \ref{thm:26.5.1} 得到。
\end{proof}

\begin{corollary}\label{cor:27.3.10}
设 $S_+$ 由 $S_1$ 生成，且 $S_1$ 有限型。则典范态射
\[
\Proj_Y(S) \to \Spec_Y(S_0)
\]
紧合。
\end{corollary}

\begin{proof}
由推论 \ref{cor:26.5.3} 得到。
\end{proof}

\subsection{射影丛}\label{sec:27.4}

\begin{construction}\label{con:27.4.1}
设 $Y$ 为概形，$E$ 为拟凝聚 $\cO_Y$-模。我们给 $\Sym_{\cO_Y}(E)$ 配上自然的 $\Z_{\ge 0}$-分次，使 $\Sym_{\cO_Y}(E)_1 \simeq E$。概形
\[
\PP_Y(E) := \Proj_Y\bigl(\Sym_{\cO_Y}(E)\bigr)
\]
称为 $Y$ 上与 $E$ 相伴的\emph{射影丛}。
\end{construction}

\begin{example}\label{ex:27.4.2}
有
\[
\PP_Y(\cO_Y^{\oplus n}) \simeq \PP^n_S := \PP^n_\Z \times S.
\]
\end{example}

\begin{yizhu}
此处 $S$ 在原书未出现定义：本小节从头到尾只以 $Y$ 为底概形，等号右边的 $\PP^n_S$ 应读作 $\PP^n_Y$，即
$\PP^n_\Z\times S$ 中的 $S$ 是 $Y$ 的笔误。紧接的下一小节另把 $S$ 定义为分次 $\cO_Y$-代数，与这里无关。
[依据：原书 p.137 例 27.4.2 与 §27.5 开头。]
\end{yizhu}

\begin{remark}\label{rem:27.4.3}
按定义，有典范态射
\[
\pi \colon V_Y(E) \setminus Y \to \PP_Y(E),
\]
其中 $Y \to V_Y(E) = \Spec_Y(\Sym_{\cO_Y}(E))$ 是对应于把 $E$ 送到零的同态 $\Sym_{\cO_Y}(E) \to \cO_Y$ 的闭浸入。
\end{remark}

\begin{exercise}\label{exc:27.4.4}
设 $L$ 是 $Y$ 上的线丛。证明：
\begin{enumerate}[label=(\arabic*),leftmargin=2.4em,itemsep=0.25ex]
  \item 典范态射 $p \colon \PP_Y(L) \to Y$ 是同构。
  \item 典范态射 $p^*L \to \cO(1)$ 是同构。
\end{enumerate}
\end{exercise}

\begin{yizhu}
构造 \ref{con:27.4.1} 中的 $E$ 不要求有限局部自由，也不要求有限型：Stacks Project 明确指出，为与 EGA II 的定义一致，射影丛只对拟凝聚的 $E$ 定义。代价是 $\PP_Y(E)$ 的纤维一般不是射影空间，故推论 \ref{cor:27.3.1} 的有限型结论必须另设「$S_1$ 有限型」这一条。
[来源：Stacks Project, Tag 01OB（Definition 27.21.1，射影丛 $\PP(E) = \underline{\Proj}_S(\Sym(E))$ 及其脚注），\url{https://stacks.math.columbia.edu/tag/01OB}（2026-11-17 访问）。]
\end{yizhu}

\subsection{映入相对 Proj 的态射}\label{sec:27.5}

本小节固定概形 $Y$ 与 $\Z_{\ge 0}$-分次拟凝聚 $\cO_Y$-代数 $S$。记 $\cO := \cO_{\Proj_Y(S)}$。

回忆：若 $S_+$ 由 $S_1$ 生成，则有重言线丛 $\cO(1)$ 与典范同态
\[
S \to \bigoplus_{n \ge 0} p_*\bigl(\cO(n)\bigr).
\]
利用伴随对 $(p^*, p_*)$，得到分次 $\cO$-同态
\begin{equation}\label{eq:27.4}
\phi \colon p^*(S) \to \bigoplus_{n \in \Z_{\ge 0}} \cO(n) \simeq \Sym_{\cO}\bigl(\cO(1)\bigr),
\end{equation}
使得
\[
\phi_1 \colon p^*(S_1) \to \cO(1)
\]
是满射（推论 \ref{cor:27.3.7}）。可以证明，数对 $(\cO(1), \phi)$ 刻画了 $\Proj_Y(S)$。

\begin{construction}\label{con:27.5.1}
对任意 $Y$-概形 $q \colon X \to Y$，记 $F_{Y,S}(X)$ 为如下数对的同构类之集：数对 $(L, \psi)$，其中 $L$ 是 $X$ 上的线丛，
\[
\psi \colon q^*(S) \to \Sym_{\cO_X}(L)
\]
是分次 $\cO_X$-同态，且 $\psi_1 \colon q^*(S_1) \to L$ 是满射。

对 $Y$-概形之间的态射 $f \colon X' \to X$，有典范映射
\[
F_{Y,S}(X) \to F_{Y,S}(X'), \qquad (L, \psi) \mapsto \bigl(f^*(L), f^*(\psi)\bigr),
\]
其中我们沿用上述记号，并用 $f^*(\psi)$ 表示下述合成：
\[
(q')^*(S) \simeq f^* q^*(S) \xrightarrow{\ f^*(\psi)\ } f^*\bigl(\Sym_{\cO_X}(L)\bigr) \simeq \Sym_{\cO_{X'}}\bigl(f^*(L)\bigr).
\]
这定义了函子
\[
F_{Y,S} \colon \mathsf{Sch}_Y^{\op} \to \Set.
\]
\end{construction}

\begin{theorem}\label{thm:27.5.2}
设 $S_+$ 由 $S_1$ 生成。则函子 $F_{Y,S}$ 被 $\Proj_Y(S)$ 表示，即有自然同构
\[
F_{Y,S} \simeq \Hom_Y\bigl(-, \Proj_Y(S)\bigr).
\]
此外，恒等态射 $\id_{\Proj_Y(S)}$ 对应元素
\[
\bigl(\cO(1), \phi\bigr) \in F_{Y,S}\bigl(\Proj_Y(S)\bigr),
\]
其中 $\phi$ 是典范同态 (27.4)。
\end{theorem}

为证明该定理，我们将构造互为逆的自然变换
\[
\Hom_Y\bigl(-, \Proj_Y(S)\bigr) \leftrightarrows F_{Y,S}.
\]

\begin{construction}\label{con:27.5.3}
设 $S_+$ 由 $S_1$ 生成。右向的自然变换由下述合成给出：
\[
\Hom_Y\bigl(-, \Proj_Y(S)\bigr) \to \Hom\bigl(F_{Y,S}(\Proj_Y(S)), F_{Y,S}(-)\bigr) \xrightarrow{\ \mathrm{ev}_{(\cO(1), \phi)}\ } F_{Y,S},
\]
其中最后一个映射是在典范元素 $(\cO(1), \phi) \in F_{Y,S}$ 处的求值。

具体地说，对 $Y$-概形 $q \colon X \to Y$ 与 $Y$-态射 $r \colon X \to \Proj_Y(S)$，上述自然变换把 $r$ 送到
\[
\bigl(r^*(\cO(1)), r^*(\phi)\bigr) \in F_{Y,S}(X),
\]
其中我们滥用记号，用 $r^*(\phi)$ 表示下述合成：
\[
q^*S \simeq r^* p^* S \xrightarrow{\ r^*(\phi)\ } r^*\Sym_{\cO}\bigl(\cO(1)\bigr) \simeq \Sym_{\cO_X}\bigl(r^*(\cO(1))\bigr).
\]
\end{construction}

\begin{construction}\label{con:27.5.4}
今定义左向的自然变换。为便于日后引用，先去掉「$S_+$ 由 $S_1$ 生成」的假设。

设 $q \colon X \to Y$ 为 $Y$-概形，$L$ 为 $X$ 上的线丛。对任意分次同态
\[
\psi \colon q^*S \to \Sym_{\cO_X}(L),
\]
有态射
\[
r_{L,\psi} \colon U_+(\psi) \to \Proj_X(q^*S) \simeq \Proj_Y(S) \times_Y X \to \Proj_Y(S),
\]
其中
\[
U_+(\psi) \subseteq \Proj_X\bigl(\Sym_{\cO_X}(L)\bigr) \simeq X
\]
视为 $X$ 的开子概形。可以验证：该开子概形与态射 $r_{L,\psi}$ 都只依赖 $(L, \psi)$ 的同构类。

当 $(L, \psi) \in F_{Y,S}(X)$，即 $\psi_1$ 是满射时，可见 $U_+(\psi) = X$。于是得到映射
\[
F_{Y,S}(X) \to \Hom_Y\bigl(X, \Proj_Y(S)\bigr), \qquad (L, \psi) \mapsto r_{L,\psi}.
\]
\end{construction}

上述构造对 $X$ 具有函子性，故得到自然变换
\[
F_{Y,S} \to \Hom_Y\bigl(-, \Proj_Y(S)\bigr).
\]

\begin{proof}[定理 \ref{thm:27.5.2} 的证明]
设 $q \colon X \to Y$ 为 $Y$-概形。只需证如上构造的映射
\[
\Hom_Y\bigl(X, \Proj_Y(S)\bigr) \rightleftarrows F_{Y,S}(X)
\]
互为逆。

注意有典范等同
\[
\Hom_Y\bigl(X, \Proj_Y(S)\bigr) \simeq \Hom_X\bigl(X, \Proj_Y(S) \times_Y X\bigr) \simeq \Hom_X\bigl(X, \Proj_X(q^*S)\bigr)
\]
以及
\[
F_{Y,S}(X) \simeq F_{X, q^*S}(X).
\]
于是可用 $(X, q^*S)$ 替换 $(Y, S)$，从而归约到 $X = Y$ 的情形。换言之，只需证
\begin{equation}\label{eq:27.5}
\Hom_Y\bigl(Y, \Proj_Y(S)\bigr) \to F_{Y,S}(Y), \qquad r \mapsto \bigl(r^*(\cO(1)), r^*(\phi)\bigr)
\end{equation}
与
\begin{equation}\label{eq:27.6}
F_{Y,S}(Y) \to \Hom_Y\bigl(Y, \Proj_Y(S)\bigr), \qquad (L, \psi) \mapsto r_{(L,\psi)}
\end{equation}
互为逆。

先证 $(27.6) \circ (27.5) = \id$。任意元素 $r \in \Hom_Y(Y, \Proj_Y(S))$ 都是闭浸入，因为 $\Proj_Y(S) \to Y$ 分离（习题 18.2.11(2)）。故只需识别满射
\[
\cO \to r_*\cO_Y \quad \text{与} \quad \cO \to r_{(r^*(\cO(1)), r^*\phi),*}\cO_Y.
\]
展开定义，第二个满射由下述合成给出：
\[
\cO \simeq \widetilde{S} \simeq \widetilde{r^* p^*(S)} \xrightarrow{\ \widetilde{r^*(\phi)}\ } \widetilde{\bigl(r^*\Sym_{\cO}(\cO(1))\bigr)}.
\]
于是该断言来自下述同构：
\[
\begin{aligned}
r_*\cO_Y &\simeq \bigl(p_*(r_*\cO_Y)_{\mathrm{gr}}\bigr)\widetilde{\ } \simeq \Bigl(\bigoplus_{n \in \Z_{\ge 0}} p_*\bigl(r_*\cO_Y \otimes_\cO \cO(n)\bigr)\Bigr)\widetilde{\ }\\
&\simeq \Bigl(\bigoplus_{n \in \Z_{\ge 0}} p_*\bigl(r_*(\cO_Y \otimes_{r^*\cO} r^*\cO(n))\bigr)\Bigr)\widetilde{\ } \simeq \bigl(r^*\Sym_{\cO}(\cO(1))\bigr)\widetilde{\ },
\end{aligned}
\]
其中
\begin{itemize}[leftmargin=2em,itemsep=0.25ex]
  \item 第一个同构由推论 \ref{cor:27.3.4} 得到；
  \item 第二个同构由 $(-)\widetilde{\ }$ 的构造得到；
  \item 第三个同构由下面的习题 \ref{exc:27.5.5} 得到；
  \item 最后一个同构是显然的。
\end{itemize}

剩下要证 $(27.5) \circ (27.6) = \id$。这由推论 \ref{cor:27.3.9}(1) 与习题 \ref{exc:27.4.4} 得到。
\end{proof}

\begin{exercise}\label{exc:27.5.5}
设 $r \colon (Y, \cO_Y) \to (Z, \cO_Z)$ 是环化空间之间的态射，$M$（相应地 $N$）为 $\cO_Y$-模（相应地 $\cO_Z$-模）。
\begin{enumerate}[label=(\arabic*),leftmargin=2.4em,itemsep=0.25ex]
  \item 构造典范的 $\cO_Z$-线性态射
  \[
  r_*M \otimes_{\cO_Z} N \to r_*\bigl(M \otimes_{\cO_Y} r^*N\bigr).
  \]
  \item 设 $N$ 是局部自由 $\cO_Z$-模。证明上述态射是同构。该同构称为（投影公式的一个特例）\emph{投影公式}。
\end{enumerate}
\end{exercise}

\begin{corollary}\label{cor:27.5.6}
设 $Y$ 为概形，$E$ 为拟凝聚 $\cO_Y$-模。对任意 $Y$-概形 $q \colon X \to Y$，下述两集之间有典范双射：
\begin{enumerate}[label=(\alph*),leftmargin=2.6em,itemsep=0.25ex]
  \item $Y$-态射 $X \to \PP_Y(E)$ 之集。
  \item 数对 $(L, \psi)$ 的同构类之集，其中 $L$ 是 $X$ 上的线丛，且
  \[
  q^*(E) \to L
  \]
  是 $\cO_X$-模之间的满射。
\end{enumerate}
\end{corollary}

\begin{yizhu}
定理 \ref{thm:27.5.2} 的证明之所以能归约到 $X = Y$，靠的是引理 \ref{lem:27.1.2}：基变换后仍是同一个相对 $\Proj$ 构造。Stacks Project 在一般情形下改用「$d$ 倍」的四元组 $(d, f \colon T \to S, \cL, \psi)$，其中 $\psi \colon f^*\mathcal A^{(d)} \to \bigoplus_{n \ge 0} \cL^{\otimes n}$ 且 $f^*\mathcal A_d \to \cL$ 满；本节的假设允许一律取 $d = 1$。
不同 $d$ 给出的函子表示同一概形，故 $d$ 只是记账上的方便选择。

[来源：Stacks Project, Tag 01NS（Section 27.16，Relative Proj as a functor：四元组的严格等价类与表示函子），\url{https://stacks.math.columbia.edu/tag/01NS}（2026-11-17 访问）。]
\end{yizhu}

\subsection{Segre 嵌入}\label{sec:27.6}

\begin{construction}\label{con:27.6.1}
设 $Y$ 为概形，$(E_i)_{i=1}^n$ 为拟凝聚 $\cO_Y$-模的有限族。对每个 $i$，记 $P_i := \PP_Y(E_i)$，并记 $p_i \colon P_i \to Y$ 为典范态射。设
\[
\varphi_i \colon p_i^*E_i \to \cO_{P_i}(1)
\]
为典范满射。考虑纤维积
\[
P := P_1 \times_Y \cdots \times_Y P_n,
\]
以及典范投影 $\mathrm{pr}_i \colon P \to P_i$ 与 $r \colon P \to Y$。设 $\varphi$ 为下述合成：
\[
\begin{aligned}
r^*(E_1 \otimes_{\cO_Y} \cdots \otimes_{\cO_Y} E_n) &\simeq \mathrm{pr}_1^*(p_1^*E_1) \otimes_{\cO_P} \cdots \otimes_{\cO_P} \mathrm{pr}_n^*(p_n^*E_n)\\
&\xrightarrow{\ \otimes\, \mathrm{pr}_i^*(\varphi_i)\ } \mathrm{pr}_1^*\bigl(\cO_{P_1}(1)\bigr) \otimes_{\cO_P} \cdots \otimes_{\cO_P} \mathrm{pr}_n^*\bigl(\cO_{P_n}(1)\bigr) =: \cL.
\end{aligned}
\]
\end{construction}

注意 $\cL$ 是 $P$ 上的线丛。

\begin{exercise}\label{exc:27.6.2}
证明：
\begin{enumerate}[label=(\arabic*),leftmargin=2.4em,itemsep=0.25ex]
  \item 态射 $\varphi$ 是满射。
  \item 数对 $(\cL, \varphi)$ 对应一个闭浸入
  \[
  \PP_Y(E_1) \times_Y \cdots \times_Y \PP_Y(E_n) \to \PP_Y(E_1 \otimes_{\cO_Y} \cdots \otimes_{\cO_Y} E_n).
  \]
\end{enumerate}
\end{exercise}

该闭浸入称为 \emph{Segre 嵌入}。

\begin{yizhu}
平凡情形 $Y = \Spec \Z$、$n = 2$ 值得先算：此时 $P_i = \PP^{n_i}$，$\cL = \mathrm{pr}_1^*\cO(1) \otimes \mathrm{pr}_2^*\cO(1)$，它由 $(n_1+1)(n_2+1)$ 个截面 $X_iY_j$ 生成，故给出闭浸入 $\PP^{n_1} \times \PP^{n_2} \to \PP^{n_1 n_2 + n_1 + n_2}$。闭性是这样看的：$\varphi^{-1}(D_+(Z_{i + (n_1+1)j})) = D_+(X_i) \times D_+(Y_j)$，故 $\varphi$ 仿射，而相应的坐标环映射 $Z_i/Z_0 \mapsto X_i/X_0$、$Z_{(n_1+1)j}/Z_0 \mapsto Y_j/Y_0$ 是满射。
换个说法：由推论 \ref{cor:27.5.6}，闭浸入等价于给出线丛与满射 $q^*(E_1 \otimes_{\cO_Y} \cdots \otimes_{\cO_Y} E_n) \to \cL$，而 Segre 嵌入只是 $E_i = \cO_{\PP^{n_i}}(1)$ 的显式版本。

[来源：Stacks Project, Tag 01WD（Lemma 27.13.6，Segre embedding），\url{https://stacks.math.columbia.edu/tag/01WD}（2026-11-17 访问）。]
\end{yizhu}
